\documentclass[10pt]{extarticle}

% ============================================================
% Pacotes
% ============================================================
\usepackage[utf8]{inputenc}
\usepackage[T1]{fontenc}
\usepackage[brazil]{babel}
\usepackage{amsmath,amssymb}
\usepackage{enumitem}
\usepackage{xcolor}
\usepackage[margin=1.5cm]{geometry}
\usepackage{titlesec}
\usepackage{fancyhdr}
\usepackage{hyperref}

% ============================================================
% Paleta de cores
% ============================================================
\definecolor{darkblue}{RGB}{0,51,102}
\definecolor{accentorange}{RGB}{230,126,34}

% ============================================================
% Estilo de títulos
% ============================================================
\titleformat{\section}
  {\normalfont\Large\bfseries\color{darkblue}}
  {\thesection}{1em}{}
\titleformat{\subsection}
  {\normalfont\large\bfseries\color{accentorange}}
  {\thesubsection}{1em}{}

% ============================================================
% Cabeçalho e rodapé
% ============================================================
\pagestyle{fancy}
\fancyhf{}
\renewcommand{\headrulewidth}{0.4pt}
\renewcommand{\footrulewidth}{0.4pt}
\fancyhead[L]{\color{darkblue}\small Problem Set 1}
\fancyhead[R]{\color{darkblue}\small PCC190 - Secure Machine Learning}
\fancyfoot[C]{\thepage}

% ============================================================
% Estilo da lista de questões
% ============================================================
\newlist{questoes}{enumerate}{1}
\setlist[questoes,1]{label=\textbf{\color{darkblue}\arabic*.},leftmargin=1.5cm,itemsep=0.8em}

\newlist{itens}{enumerate}{1}
\setlist[itens,1]{label=\alph*),leftmargin=1cm}

\begin{document}

% ============================================================
% Cabeçalho do documento
% ============================================================
\begin{center}
    {\LARGE\bfseries\color{darkblue}Problem Set 1}\\[0.3cm]
    {\large\color{accentorange} PCC190 - Secure Machine Learning}\\[0.5cm]
\end{center}


% ============================================================
% Questões
% ============================================================
\begin{questoes}

\item Why is machine learning important nowadays? What is its concomitant risk?

\item How does the book define the problem of secure machine learning? Are there other possible definitions? What are they?

\item What are the key strengths of learning approaches, and how can they be subverted?

\item What is the main reason machine learning techniques are being applied to a growing number of systems?

\item What is a malfeasance detection problem?

\item Give examples of malfeasance detection problems that can be solved with machine learning techniques, and explain how those techniques solve them.

\item According to the book, what is the primary vulnerability in learners that attackers can exploit?

\item What are the two fundamental questions regarding machine learning security?

\item Describe how a spam filter based on machine learning works.

\item Describe the operation of a network anomaly detector based on PCA. How can an attacker evade such a system?

\item Why would privacy be a concern in machine learning systems? Give examples.

\item What are the steps involved in the evaluation of a computer system's security?

\item What is computer security concerned about?

\item In your own words, describe the following classical security principles:
\begin{itens}
    \item Proactive analysis
    \item Kerckhoffs's Principle
    \item Conservative design
    \item Threat modeling (how is a threat model constructed?)
\end{itens}

\end{questoes}

\end{document}